\documentclass[10pt,a4paper]{amsart}
\usepackage{amsmath,amssymb,amsthm,graphicx}
\usepackage{verbatim}
\newenvironment{reference}{\comment}{\endcomment}
%\newenvironment{reference}{}{}
\newenvironment{slogan}{\comment}{\endcomment}
\newenvironment{history}{\comment}{\endcomment}

% For commutative diagrams we use Xy-pic
\usepackage[all]{xy}

% We use 2cell for 2-commutative diagrams.
\xyoption{2cell}
% Theorem environments.
%
\theoremstyle{plain}
\newtheorem{theorem}[subsection]{Theorem}
\newtheorem{proposition}[subsection]{Proposition}
\newtheorem{lemma}[subsection]{Lemma}

\theoremstyle{definition}
\newtheorem{definition}[subsection]{Definition}
\newtheorem{example}[subsection]{Example}
\newtheorem{exercise}[subsection]{Exercise}
\newtheorem{situation}[subsection]{Situation}

\theoremstyle{remark}
\newtheorem{remark}[subsection]{Remark}
\newtheorem{remarks}[subsection]{Remarks}

\numberwithin{equation}{subsection}

% Macros
%
\def\lim{\mathop{\mathrm{lim}}\nolimits}
\def\colim{\mathop{\mathrm{colim}}\nolimits}
\def\Spec{\mathop{\mathrm{Spec}}}
\def\Hom{\mathop{\mathrm{Hom}}\nolimits}
\def\Ext{\mathop{\mathrm{Ext}}\nolimits}
\def\SheafHom{\mathop{\mathcal{H}\!\mathit{om}}\nolimits}
\def\SheafExt{\mathop{\mathcal{E}\!\mathit{xt}}\nolimits}
\def\Sch{\mathit{Sch}}
\def\Mor{\mathop{\mathrm{Mor}}\nolimits}
\def\Ob{\mathop{\mathrm{Ob}}\nolimits}
\def\Sh{\mathop{\mathit{Sh}}\nolimits}
\def\NL{\mathop{N\!L}\nolimits}
\def\CH{\mathop{\mathrm{CH}}\nolimits}
\def\proetale{{pro\text{-}\acute{e}tale}}
\def\etale{{\acute{e}tale}}
\def\QCoh{\mathit{QCoh}}
\def\Ker{\mathop{\mathrm{Ker}}}
\def\Im{\mathop{\mathrm{Im}}}
\def\Coker{\mathop{\mathrm{Coker}}}
\def\Coim{\mathop{\mathrm{Coim}}}

% Boxtimes
%
\DeclareMathSymbol{\boxtimes}{\mathbin}{AMSa}{"02}

%
% Macros for moduli stacks/spaces
%
\def\QCohstack{\mathcal{QC}\!\mathit{oh}}
\def\Cohstack{\mathcal{C}\!\mathit{oh}}
\def\Spacesstack{\mathcal{S}\!\mathit{paces}}
\def\Quotfunctor{\mathrm{Quot}}
\def\Hilbfunctor{\mathrm{Hilb}}
\def\Curvesstack{\mathcal{C}\!\mathit{urves}}
\def\Polarizedstack{\mathcal{P}\!\mathit{olarized}}
\def\Complexesstack{\mathcal{C}\!\mathit{omplexes}}
% \Pic is the operator that assigns to X its picard group, usage \Pic(X)
% \Picardstack_{X/B} denotes the Picard stack of X over B
% \Picardfunctor_{X/B} denotes the Picard functor of X over B
\def\Pic{\mathop{\mathrm{Pic}}\nolimits}
\def\Picardstack{\mathcal{P}\!\mathit{ic}}
\def\Picardfunctor{\mathrm{Pic}}
\def\Deformationcategory{\mathcal{D}\!\mathit{ef}}
\title{Three equivalent characterizations of a locally ringed site}
\author{The Stacks Project Authors}
\date{stacks-project:sites-modules.tex:local-snapshot-20261012; GNU FDL 1.2 or later}
\begin{document}
\maketitle
\noindent\textit{Editorial note.} All delivered mathematical text is the author's own, copied verbatim from the pinned Stacks Project chapter file sites-modules.tex; no statement or proof step was rewritten, reordered or removed.\par\medskip
\noindent\textit{Editorial note.} The counted span contains no external \textbackslash{}ref, so no navigation stubs were needed.\par\medskip
\noindent\textit{Editorial note.} This is the first card this lane authors from the source stacks-sites-modules; its preamble is taken as two verbatim slices of the chapter's own preamble.tex (lines 8-18 and 154-217), which include the verbatim/reference/slogan/history environment definitions required by the chapter body.\par\medskip
\noindent\textit{Allowed standard prerequisites (editorial).}\begin{itemize}
\item Sheaves on a site and coverings. Stacks Project, Sites, Definition of a site and of a sheaf (coverings and the sheaf condition). Supplies the language of coverings in which all three formulations are expressed.
\end{itemize}\medskip
\begin{lemma}
\label{lemma-locally-ringed}
Let $(\mathcal{C}, \mathcal{O})$ be a ringed site. The following
are equivalent
\begin{enumerate}
\item For every object $U$ of $\mathcal{C}$ and $f \in \mathcal{O}(U)$
there exists a covering $\{U_j \to U\}$ such that for each $j$
either $f|_{U_j}$ is invertible or $(1 - f)|_{U_j}$ is invertible.
\item For $U \in \Ob(\mathcal{C})$, $n \geq 1$, and
$f_1, \ldots, f_n \in \mathcal{O}(U)$ which generate the unit ideal
in $\mathcal{O}(U)$ there exists a covering $\{U_j \to U\}$
such that for each $j$ there exists an $i$ such that $f_i|_{U_j}$
is invertible.
\item The map of sheaves of sets
$$
(\mathcal{O} \times \mathcal{O})
\amalg
(\mathcal{O} \times \mathcal{O})
\longrightarrow
\mathcal{O} \times \mathcal{O}
$$
which maps $(f, a)$ in the first component to $(f, af)$ and
$(f, b)$ in the second component to $(f, b(1 - f))$ is surjective.
\end{enumerate}
\end{lemma}

\begin{proof}
It is clear that (2) implies (1). To show that (1) implies (2) we argue by
induction on $n$. The first case is $n = 2$ (since $n = 1$ is trivial).
In this case we have $a_1f_1 + a_2f_2 = 1$ for some
$a_1, a_2 \in \mathcal{O}(U)$. By assumption we can find a covering
$\{U_j \to U\}$ such that for each $j$
either $a_1f_1|_{U_j}$ is invertible or $a_2f_2|_{U_j}$ is invertible.
Hence either $f_1|_{U_j}$ is invertible or $f_2|_{U_j}$ is invertible
as desired. For $n > 2$ we have
$a_1f_1 + \ldots + a_nf_n = 1$ for some $a_1, \ldots, a_n \in \mathcal{O}(U)$.
By the case $n = 2$ we see that we have some covering $\{U_j \to U\}_{j \in J}$
such that for each $j$ either $f_n|_{U_j}$ is invertible or
$a_1f_1 + \ldots + a_{n - 1}f_{n - 1}|_{U_j}$ is invertible.
Say the first case happens for $j \in J_n$. Set $J' = J \setminus J_n$.
By induction hypothesis, for each $j \in J'$ we can find a covering
$\{U_{jk} \to U_j\}_{k \in K_j}$ such that for each $k \in K_j$ there
exists an $i \in \{1, \ldots, n - 1\}$ such that
$f_i|_{U_{jk}}$ is invertible. By the axioms of a site the family of
morphisms
$\{U_j \to U\}_{j \in J_n} \cup \{U_{jk} \to U\}_{j \in J', k \in K_j}$
is a covering which has the desired property.

\medskip\noindent
Assume (1). To see that the map in (3) is surjective, let
$(f, c)$ be a section of $\mathcal{O} \times \mathcal{O}$ over $U$.
By assumption there exists a covering $\{U_j \to U\}$ such that
for each $j$ either $f$ or $1 - f$ restricts to an invertible section.
In the first case we can take $a = c|_{U_j} (f|_{U_j})^{-1}$, and
in the second case we can take $b = c|_{U_j} (1 - f|_{U_j})^{-1}$.
Hence $(f, c)$ is in the image of the map on each of the members.
Conversely, assume (3) holds. For any $U$ and $f \in \mathcal{O}(U)$
there exists a covering $\{U_j \to U\}$ of $U$ such that the
section $(f, 1)|_{U_j}$ is in the image of the map in (3) on sections
over $U_j$. This means precisely that either $f$ or $1 - f$ restricts
to an invertible section over $U_j$, and we see that (1) holds.
\end{proof}

\end{document}
